\section{Temporización y medición de tiempo}
\label{sec:temporizacion}

\subsection{Frecuencia de reloj y conversión de ciclos a décimas}

La Go Board tiene un oscilador de 25~MHz conectado al pin 15 de la FPGA, y el SoC usa
ese reloj sin dividirlo. Todas las cantidades de tiempo del juego se derivan de ese
valor:

\begin{align*}
  f_{\text{clk}} &= 25\,000\,000~\text{Hz}, \qquad T_{\text{clk}} = \frac{1}{f_{\text{clk}}} = 40~\text{ns} \\[4pt]
  N_{\text{décima}} &= f_{\text{clk}} \cdot 0{,}1~\text{s} = 2\,500\,000~\text{ciclos} \\[4pt]
  N_{3\,\text{s}} &= f_{\text{clk}} \cdot 3~\text{s} = 75\,000\,000~\text{ciclos}
\end{align*}

Si el contador marca $c_0$ cuando se enciende el LED objetivo y $c_1$ cuando se detecta
el botón, el tiempo de respuesta en décimas de segundo es

\[
  t = \left\lfloor \frac{(c_1 - c_0) \bmod 2^{32}}{2\,500\,000} \right\rfloor .
\]

\begin{itemize}
  \item \textbf{La resta se hace sin signo, en 32 bits} (\code{sub}, que da el
        resultado módulo $2^{32}$). Si el contador dio la vuelta entre $c_0$ y $c_1$,
        la resta sigue dando la cantidad correcta de ciclos. El único caso que falla
        es una espera de más de $2^{32}$ ciclos, unos 171,8 segundos.
  \item \textbf{La división trunca}, como un cronómetro: 0,29~s se muestra como
        \code{02}. Así un valor mostrado de $n$ décimas significa que el tiempo real
        está entre $n$ y $n+1$ décimas.
  \item \textbf{La división se hace por restas sucesivas} (subrutina \code{divu}),
        porque el procesador no tiene instrucción de división. Cada vuelta del bucle
        toma 9 ciclos, así que convertir 9,9~s (99 vueltas) toma unos 900 ciclos,
        36~$\mu$s. Esa cuenta ocurre después de leer $c_1$, así que no afecta la
        medición.
\end{itemize}

Las constantes están con nombre al comienzo de \code{game.s}
(código~\ref{lst:constantes}). Cambiar la frecuencia del reloj solo exige cambiar
\code{WAIT\_CYCLES} y \code{TENTH\_CYCLES}.

\begin{lstlisting}[language=rvasm, caption={Constantes de tiempo en \code{sw/game.s}.}, label={lst:constantes}]
.equ ROUNDS,          10          # respuestas correctas por partida
.equ WAIT_CYCLES,     75000000    # fase con los 4 LEDs: 3 s = 3 * 25_000_000
.equ TENTH_CYCLES,    2500000     # ciclos en una decima de segundo = 25_000_000 / 10
.equ DEBOUNCE_CYCLES, 500000      # 20 ms con los botones sueltos antes de mostrar el objetivo
.equ BLINK_CYCLES,    6250000     # 0,25 s: cada cuanto cambian los LEDs mientras se espera
\end{lstlisting}

\subsection{Espera de tres segundos}

La subrutina \code{wait\_cycles} (código~\ref{lst:wait}) lee el contador al entrar y
vuelve cuando la diferencia con la lectura actual alcanza el valor pedido en
\code{a0}. La comparación \code{bltu} es sin signo, por la misma razón que la resta.

\begin{lstlisting}[language=rvasm, caption={Espera de \code{a0} ciclos.}, label={lst:wait}]
wait_cycles:
    lw   t0, CYCLES(s0)             # t0 = inicio
wc_loop:
    lw   t1, CYCLES(s0)
    sub  t1, t1, t0                 # t1 = ciclos transcurridos
    bltu t1, a0, wc_loop
    ret
\end{lstlisting}

Cada vuelta del bucle toma 8 ciclos (\code{lw} 3, \code{sub} 2, \code{bltu} tomado 3),
así que la espera termina a lo más 8 ciclos, 320~ns, después de cumplirse los
75.000.000. Tras ella el programa exige además 20~ms seguidos sin ningún botón
apretado (siguiente subsección). La fase con los cuatro LEDs dura entonces
$3{,}02$~s si nadie toca los botones, y más si alguien los mantiene apretados, lo que
cumple el ``al menos tres segundos'' del enunciado.

\subsection{Rebotes de los botones}

Un botón mecánico rebota: al apretarlo o soltarlo, el contacto se abre y cierra varias
veces durante algunos milisegundos. El juego solo necesita el primer contacto para
detener el cronómetro, así que los rebotes al apretar no afectan la medición. Los
rebotes al soltar sí son un problema. Si el jugador suelta el botón justo cuando
aparece el LED objetivo, un rebote se leería como una pulsación nueva, con un tiempo de
casi cero.

La subrutina \code{wait\_release} se ejecuta antes de encender el LED objetivo. Vuelve
solo cuando los botones llevan \code{DEBOUNCE\_CYCLES} (20~ms) seguidos sin apretarse:
cada lectura con un botón apretado reinicia la cuenta. Así se filtran los rebotes y,
además, no se puede ``hacer trampa'' dejando un botón apretado antes de que aparezca el
objetivo.

\subsection{Medición del tiempo de respuesta}

El código~\ref{lst:medicion} muestra el núcleo de la ronda. El cronómetro se inicia
con la instrucción siguiente a la que enciende el LED objetivo, y se detiene con la
instrucción siguiente a la que detecta el botón.

\begin{lstlisting}[language=rvasm, caption={Encendido del objetivo y medición.}, label={lst:medicion}]
    sw   a3, LEDS(s0)               #    solo el objetivo encendido
    lw   a2, CYCLES(s0)             # 3) y en ese mismo instante parte el cronometro
wait_press:
    lw   t0, BUTTONS(s0)
    beqz t0, wait_press             # 4) esperar un boton
    lw   t1, CYCLES(s0)             #    fin del cronometro
    bne  t0, a3, wrong              #    boton incorrecto (o mas de uno a la vez)
    sub  a0, t1, a2                 # 5) ciclos de reaccion
\end{lstlisting}

El error de la medición se acota con la tabla~\ref{tab:cpi}:

\begin{itemize}
  \item Entre que se escribe el LED y se lee $c_0$ pasan del orden de 2 ciclos.
  \item El bucle \code{wait\_press} da una vuelta cada 6 ciclos (\code{lw} 3 y
        \code{beqz} tomado 3), así que el botón se detecta a lo más 6 ciclos después de
        apretarse.
  \item Entre la detección y la lectura de $c_1$ pasan unos 4 ciclos.
\end{itemize}

En total el error es de unos 10 ciclos, 0,4~$\mu$s, doscientas mil veces menor que la
resolución de una décima. La medición es, en la práctica, exacta al ciclo que la
precede.

\subsection{Promedio}
\label{sec:promedio}

El programa acumula en el registro \code{a4} las décimas de cada ronda correcta. Tras
\code{ROUNDS} respuestas correctas, el promedio se calcula con la misma subrutina de
división:

\[
  \bar{t} = \left\lfloor \frac{\sum_{i=1}^{\text{ROUNDS}} t_i}{\text{ROUNDS}} \right\rfloor~\text{décimas}.
\]

Se promedian los tiempos en décimas y no los ciclos porque así el promedio es el de los
diez números que el jugador vio en los displays. Además, la suma de diez tiempos en
ciclos puede desbordar 32 bits si las rondas son lentas: $2^{32}$ ciclos son 171,8~s en
total. La suma en décimas no tiene ese problema.

Los displays tienen dos dígitos, así que un tiempo de 10~s o más se muestra como
\code{99}. Esa saturación es solo de la visualización: a la suma entra el valor
verdadero. Por ejemplo, en la simulación de la sección~\ref{sec:verificacion}, una
ronda de 12,0~s se muestra como \code{99} pero suma 120 décimas al promedio.

Las rondas con botón incorrecto no modifican ni \code{a4} ni el contador de rondas
\code{a5}: la ronda se repite y solo el intento correcto aporta su tiempo.
